% Part: first-order-logic
% Chapter: syntax-and-semantics
% Section: assignments

\documentclass[../../../include/open-logic-section]{subfiles}

\begin{document}

\olfileid{fol}{syn}{ass}

\olsection{إسنادات المتغيرات}

\begin{explain}
إسناد ال!!{variable}s~${\OLId{latin-ordinary}{s}}$ يعيّن قيمة \emph{لكل} متغير، مع أن
المتغيرات لا تتناهى. وليس هذا الشمول لازمًا لذات المسألة؛ وإنما
نفرض على إسنادات ال!!{variable}s أن تشمل جميع ال!!{variable}s
تيسيرًا للتعريف والبرهان. فإن قيمة الحد~${\OLId{latin-ordinary}{t}}$، وإرضاء
!!a{formula}~$!A$ أو عدم إرضائها في !!a{structure} بحسب~${\OLId{latin-ordinary}{s}}$،
لا يتوقفان إلا على ما يعيّنه~${\OLId{latin-ordinary}{s}}$ لل!!{variable}s الواردة في~${\OLId{latin-ordinary}{t}}$
ولل!!{variable}s الحرة في~$!A$. وهذا ما نثبته في القضيتين الآتيتين.
وضبط معنى التوقف أن الإسنادين إذا اتفقا على متغيرات~${\OLId{latin-ordinary}{t}}$ كلها
أعطيا الحد قيمة واحدة، وإذا اتفقا على المتغيرات الحرة في~$!A$
كانت $!A$ مُرضاة بأحدهما إذا وفقط إذا كانت مُرضاة بالآخر.
\end{explain}

\begin{prop}\ollabel{prop:valindep}
إذا كانت ال!!{variable}s في حد~${\OLId{latin-ordinary}{t}}$ من بين ${\OLId{latin-ordinary}{x}}_{\OLMathNumeral{1}}$، و\dots،~${\OLId{latin-ordinary}{x}}_{\OLId{latin-ordinary}{n}}$،
وكان ${\OLId{latin-ordinary}{s}}_{\OLMathNumeral{1}}({\OLId{latin-ordinary}{x}}_{\OLId{latin-ordinary}{i}}) = {\OLId{latin-ordinary}{s}}_{\OLMathNumeral{2}}({\OLId{latin-ordinary}{x}}_{\OLId{latin-ordinary}{i}})$ عندما ${\OLId{latin-ordinary}{i}} = {\OLMathNumeral{1}}$، و\dots،~${\OLId{latin-ordinary}{n}}$، فإن
$\Value{{\OLId{latin-ordinary}{t}}}{{\OLId{latin-looped}{M}}}[{\OLId{latin-ordinary}{s}}_{\OLMathNumeral{1}}] = \Value{{\OLId{latin-ordinary}{t}}}{{\OLId{latin-looped}{M}}}[{\OLId{latin-ordinary}{s}}_{\OLMathNumeral{2}}]$.
\end{prop}

\begin{proof}
نبرهن بالاستقراء على تعقيد~${\OLId{latin-ordinary}{t}}$. فأما الأساس، فلا يكون ${\OLId{latin-ordinary}{t}}$ إلا
!!a{constant} أو واحدًا من المتغيرات~${\OLId{latin-ordinary}{x}}_{\OLMathNumeral{1}}$، و\dots،~${\OLId{latin-ordinary}{x}}_{\OLId{latin-ordinary}{n}}$.
فإن كان ${\OLId{latin-ordinary}{t}} = {\OLId{latin-ordinary}{c}}$، حصل
$\Value{{\OLId{latin-ordinary}{t}}}{{\OLId{latin-looped}{M}}}[{\OLId{latin-ordinary}{s}}_{\OLMathNumeral{1}}] = \Assign{{\OLId{latin-ordinary}{c}}}{{\OLId{latin-looped}{M}}} = \Value{{\OLId{latin-ordinary}{t}}}{{\OLId{latin-looped}{M}}}[{\OLId{latin-ordinary}{s}}_{\OLMathNumeral{2}}]$.
وإن كان ${\OLId{latin-ordinary}{t}} = {\OLId{latin-ordinary}{x}}_{\OLId{latin-ordinary}{i}}$، فـ${\OLId{latin-ordinary}{s}}_{\OLMathNumeral{1}}({\OLId{latin-ordinary}{x}}_{\OLId{latin-ordinary}{i}}) = {\OLId{latin-ordinary}{s}}_{\OLMathNumeral{2}}({\OLId{latin-ordinary}{x}}_{\OLId{latin-ordinary}{i}})$ بالفرض، ولذلك
$\Value{{\OLId{latin-ordinary}{t}}}{{\OLId{latin-looped}{M}}}[{\OLId{latin-ordinary}{s}}_{\OLMathNumeral{1}}] = {\OLId{latin-ordinary}{s}}_{\OLMathNumeral{1}}({\OLId{latin-ordinary}{x}}_{\OLId{latin-ordinary}{i}}) = {\OLId{latin-ordinary}{s}}_{\OLMathNumeral{2}}({\OLId{latin-ordinary}{x}}_{\OLId{latin-ordinary}{i}}) = \Value{{\OLId{latin-ordinary}{t}}}{{\OLId{latin-looped}{M}}}[{\OLId{latin-ordinary}{s}}_{\OLMathNumeral{2}}]$.

وأما خطوة الاستقراء، فنفرض
${\OLId{latin-ordinary}{t}} = \Atom{f}{t_1, \dots, t_k}$، ونفرض المطلوب ثابتًا لـ${\OLId{latin-ordinary}{t}}_{\OLMathNumeral{1}}$،
و\dots،~${\OLId{latin-ordinary}{t}}_{\OLId{latin-ordinary}{k}}$. فيكون
\begin{align*}
  \Value{{\OLId{latin-ordinary}{t}}}{{\OLId{latin-looped}{M}}}[{\OLId{latin-ordinary}{s}}_{\OLMathNumeral{1}}] & = \Value{\Atom{f}{t_1, \dots, t_k}}{{\OLId{latin-looped}{M}}}[{\OLId{latin-ordinary}{s}}_{\OLMathNumeral{1}}] \\
  & = \Assign{{\OLId{latin-ordinary}{f}}}{{\OLId{latin-looped}{M}}}(\Value{{\OLId{latin-ordinary}{t}}_{\OLMathNumeral{1}}}{{\OLId{latin-looped}{M}}}[{\OLId{latin-ordinary}{s}}_{\OLMathNumeral{1}}], \dots, \Value{{\OLId{latin-ordinary}{t}}_{\OLId{latin-ordinary}{k}}}{{\OLId{latin-looped}{M}}}[{\OLId{latin-ordinary}{s}}_{\OLMathNumeral{1}}]).
\intertext{عندما ${\OLId{latin-ordinary}{j}} = {\OLMathNumeral{1}}$، و\dots،~${\OLId{latin-ordinary}{k}}$، تكون ال!!{variable}s
    في~${\OLId{latin-ordinary}{t}}_{\OLId{latin-ordinary}{j}}$ من بين ${\OLId{latin-ordinary}{x}}_{\OLMathNumeral{1}}$، و\dots،~${\OLId{latin-ordinary}{x}}_{\OLId{latin-ordinary}{n}}$. وبفرضية الاستقراء،
    $\Value{{\OLId{latin-ordinary}{t}}_{\OLId{latin-ordinary}{j}}}{{\OLId{latin-looped}{M}}}[{\OLId{latin-ordinary}{s}}_{\OLMathNumeral{1}}] = \Value{{\OLId{latin-ordinary}{t}}_{\OLId{latin-ordinary}{j}}}{{\OLId{latin-looped}{M}}}[{\OLId{latin-ordinary}{s}}_{\OLMathNumeral{2}}]$. ومن ثم}
\Value{{\OLId{latin-ordinary}{t}}}{{\OLId{latin-looped}{M}}}[{\OLId{latin-ordinary}{s}}_{\OLMathNumeral{1}}] & = \Value{\Atom{f}{t_1, \dots, t_k}}{{\OLId{latin-looped}{M}}}[{\OLId{latin-ordinary}{s}}_{\OLMathNumeral{1}}] \\
  & = \Assign{{\OLId{latin-ordinary}{f}}}{{\OLId{latin-looped}{M}}}(\Value{{\OLId{latin-ordinary}{t}}_{\OLMathNumeral{1}}}{{\OLId{latin-looped}{M}}}[{\OLId{latin-ordinary}{s}}_{\OLMathNumeral{1}}], \dots, \Value{{\OLId{latin-ordinary}{t}}_{\OLId{latin-ordinary}{k}}}{{\OLId{latin-looped}{M}}}[{\OLId{latin-ordinary}{s}}_{\OLMathNumeral{1}}]) \\
  & = \Assign{{\OLId{latin-ordinary}{f}}}{{\OLId{latin-looped}{M}}}(\Value{{\OLId{latin-ordinary}{t}}_{\OLMathNumeral{1}}}{{\OLId{latin-looped}{M}}}[{\OLId{latin-ordinary}{s}}_{\OLMathNumeral{2}}], \dots, \Value{{\OLId{latin-ordinary}{t}}_{\OLId{latin-ordinary}{k}}}{{\OLId{latin-looped}{M}}}[{\OLId{latin-ordinary}{s}}_{\OLMathNumeral{2}}]) \\
  & = \Value{\Atom{f}{t_1, \dots, t_k}}{{\OLId{latin-looped}{M}}}[{\OLId{latin-ordinary}{s}}_{\OLMathNumeral{2}}] = \Value{{\OLId{latin-ordinary}{t}}}{{\OLId{latin-looped}{M}}}[{\OLId{latin-ordinary}{s}}_{\OLMathNumeral{2}}].
\end{align*}
\end{proof}

\begin{prop}\ollabel{prop:satindep}
إذا كانت ال!!{variable}s الحرة في~$!A$ من بين ${\OLId{latin-ordinary}{x}}_{\OLMathNumeral{1}}$، و\dots،~${\OLId{latin-ordinary}{x}}_{\OLId{latin-ordinary}{n}}$،
وكان ${\OLId{latin-ordinary}{s}}_{\OLMathNumeral{1}}({\OLId{latin-ordinary}{x}}_{\OLId{latin-ordinary}{i}}) = {\OLId{latin-ordinary}{s}}_{\OLMathNumeral{2}}({\OLId{latin-ordinary}{x}}_{\OLId{latin-ordinary}{i}})$ عندما ${\OLId{latin-ordinary}{i}} = {\OLMathNumeral{1}}$، و\dots،~${\OLId{latin-ordinary}{n}}$، فإن
$\Sat{{\OLId{latin-looped}{M}}}{!A}[{\OLId{latin-ordinary}{s}}_{\OLMathNumeral{1}}]$ إذا وفقط إذا كان $\Sat{{\OLId{latin-looped}{M}}}{!A}[{\OLId{latin-ordinary}{s}}_{\OLMathNumeral{2}}]$.
\end{prop}

\begin{proof}
والبرهان أيضًا بالاستقراء على تعقيد~$!A$. ففي الأساس تكون $!A$
ذرية، فلا تخرج $!A$ عن الصور الآتية:
\iftag{prvTrue}{$\ltrue$،}{}
\iftag{prvFalse}{$\lfalse$،}{}
$\Atom{R}{t_1, \dots, t_k}$ لمحمول ذي ${\OLId{latin-ordinary}{k}}$ مواضع~${\OLId{latin-looped}{R}}$
وحدود ${\OLId{latin-ordinary}{t}}_{\OLMathNumeral{1}}$، و\dots،~${\OLId{latin-ordinary}{t}}_{\OLId{latin-ordinary}{k}}$، أو $\eq[{\OLId{latin-ordinary}{t}}_{\OLMathNumeral{1}}][{\OLId{latin-ordinary}{t}}_{\OLMathNumeral{2}}]$ لحدين ${\OLId{latin-ordinary}{t}}_{\OLMathNumeral{1}}$
و${\OLId{latin-ordinary}{t}}_{\OLMathNumeral{2}}$. ويكفينا في الحالتين الأخيرتين إثبات جهة الذهاب من
ال!!{biconditional}؛ فإن جهة الإياب تثبت بتبادل الإسنادين.

\begin{enumerate}
\tagitem{prvTrue}{%
  \indcase{!A}{\ltrue}{كل من $\Sat{{\OLId{latin-looped}{M}}}{\indfrm}[{\OLId{latin-ordinary}{s}}_{\OLMathNumeral{1}}]$
    و$\Sat{{\OLId{latin-looped}{M}}}{\indfrm}[{\OLId{latin-ordinary}{s}}_{\OLMathNumeral{2}}]$.}}{}

\tagitem{prvFalse}{%
  \indcase{!A}{\lfalse}{كل من $\Sat/{{\OLId{latin-looped}{M}}}{!A}[{\OLId{latin-ordinary}{s}}_{\OLMathNumeral{1}}]$
    و$\Sat/{{\OLId{latin-looped}{M}}}{!A}[{\OLId{latin-ordinary}{s}}_{\OLMathNumeral{2}}]$.}}{}

\item
  \indcase{!A}{\Atom{R}{t_1, \ldots, t_k}}{افترض
    $\Sat{{\OLId{latin-looped}{M}}}{\indfrm}[{\OLId{latin-ordinary}{s}}_{\OLMathNumeral{1}}]$. عندئذ
    \[
    \langle \Value{{\OLId{latin-ordinary}{t}}_{\OLMathNumeral{1}}}{{\OLId{latin-looped}{M}}}[{\OLId{latin-ordinary}{s}}_{\OLMathNumeral{1}}], \ldots, \Value{{\OLId{latin-ordinary}{t}}_{\OLId{latin-ordinary}{k}}}{{\OLId{latin-looped}{M}}}[{\OLId{latin-ordinary}{s}}_{\OLMathNumeral{1}}] \rangle
    \in \Assign{{\OLId{latin-looped}{R}}}{{\OLId{latin-looped}{M}}}.
    \]
    عندما ${\OLId{latin-ordinary}{i}} = {\OLMathNumeral{1}}$، و\dots،~${\OLId{latin-ordinary}{k}}$، يكون $\Value{{\OLId{latin-ordinary}{t}}_{\OLId{latin-ordinary}{i}}}{{\OLId{latin-looped}{M}}}[{\OLId{latin-ordinary}{s}}_{\OLMathNumeral{1}}] =
    \Value{{\OLId{latin-ordinary}{t}}_{\OLId{latin-ordinary}{i}}}{{\OLId{latin-looped}{M}}}[{\OLId{latin-ordinary}{s}}_{\OLMathNumeral{2}}]$ بحسب \olref{prop:valindep}. ومن ثم لدينا أيضًا
    $\langle \Value{{\OLId{latin-ordinary}{t}}_{\OLMathNumeral{1}}}{{\OLId{latin-looped}{M}}}[{\OLId{latin-ordinary}{s}}_{\OLMathNumeral{2}}], \ldots, \Value{{\OLId{latin-ordinary}{t}}_{\OLId{latin-ordinary}{k}}}{{\OLId{latin-looped}{M}}}[{\OLId{latin-ordinary}{s}}_{\OLMathNumeral{2}}] \rangle
    \in \Assign{{\OLId{latin-looped}{R}}}{{\OLId{latin-looped}{M}}}$، وبالتالي $\Sat{{\OLId{latin-looped}{M}}}{\indfrm}[{\OLId{latin-ordinary}{s}}_{\OLMathNumeral{2}}]$.}
\item
  \indcase{!A}{\eq[t_1][t_2]}{افترض $\Sat{{\OLId{latin-looped}{M}}}{\indfrm}[{\OLId{latin-ordinary}{s}}_{\OLMathNumeral{1}}]$.
    عندئذ $\Value{{\OLId{latin-ordinary}{t}}_{\OLMathNumeral{1}}}{{\OLId{latin-looped}{M}}}[{\OLId{latin-ordinary}{s}}_{\OLMathNumeral{1}}] = \Value{{\OLId{latin-ordinary}{t}}_{\OLMathNumeral{2}}}{{\OLId{latin-looped}{M}}}[{\OLId{latin-ordinary}{s}}_{\OLMathNumeral{1}}]$. ومن ثم
    \begin{align*}
      \Value{{\OLId{latin-ordinary}{t}}_{\OLMathNumeral{1}}}{{\OLId{latin-looped}{M}}}[{\OLId{latin-ordinary}{s}}_{\OLMathNumeral{2}}] & = \Value{{\OLId{latin-ordinary}{t}}_{\OLMathNumeral{1}}}{{\OLId{latin-looped}{M}}}[{\OLId{latin-ordinary}{s}}_{\OLMathNumeral{1}}]
      & \text{(بحسب \olref{prop:valindep})} \\
      & = \Value{{\OLId{latin-ordinary}{t}}_{\OLMathNumeral{2}}}{{\OLId{latin-looped}{M}}}[{\OLId{latin-ordinary}{s}}_{\OLMathNumeral{1}}]
      & \text{(لأن $\Sat{{\OLId{latin-looped}{M}}}{\eq[{\OLId{latin-ordinary}{t}}_{\OLMathNumeral{1}}][{\OLId{latin-ordinary}{t}}_{\OLMathNumeral{2}}]}[{\OLId{latin-ordinary}{s}}_{\OLMathNumeral{1}}]$)}\\
      &= \Value{{\OLId{latin-ordinary}{t}}_{\OLMathNumeral{2}}}{{\OLId{latin-looped}{M}}}[{\OLId{latin-ordinary}{s}}_{\OLMathNumeral{2}}]
      & \text{(بحسب \olref{prop:valindep})،}
    \end{align*}
    ولذلك $\Sat{{\OLId{latin-looped}{M}}}{\eq[{\OLId{latin-ordinary}{t}}_{\OLMathNumeral{1}}][{\OLId{latin-ordinary}{t}}_{\OLMathNumeral{2}}]}[{\OLId{latin-ordinary}{s}}_{\OLMathNumeral{2}}]$.}
\end{enumerate}

ونفرض في خطوة الاستقراء أن $\Sat{{\OLId{latin-looped}{M}}}{!B}[{\OLId{latin-ordinary}{s}}_{\OLMathNumeral{1}}]$ إذا وفقط إذا كان
$\Sat{{\OLId{latin-looped}{M}}}{!B}[{\OLId{latin-ordinary}{s}}_{\OLMathNumeral{2}}]$ لكل ال!!{formula}s~$!B$ الأقل تعقيدًا من~$!A$،
كلما اتفق الإسنادان على متغيراتها الحرة. ثم نفصّل الحالات بحسب
المؤثر الرئيس لـ~$!A$. ويكفي في كل حالة إثبات جهة واحدة من
ال!!{biconditional}، إذ تثبت الأخرى بتبادل الإسنادين. وفي غير
حالتي المكممين نطبّق الفرض الاستقرائي على ال!!{subformula}s
~$!B$ من~$!A$. فكل متغير حر في~$!B$ حر في~$!A$؛ ومن ثم، إذا
اتفق ${\OLId{latin-ordinary}{s}}_{\OLMathNumeral{1}}$ و${\OLId{latin-ordinary}{s}}_{\OLMathNumeral{2}}$ على المتغيرات الحرة في~$!A$، اتفقا على المتغيرات
الحرة في~$!B$، فجاز تطبيق الفرض الاستقرائي على~$!B$.

\begin{enumerate}
\tagitem{defNot}{}{%
  \iftag{probNot}{%
    \indcase!{!A}{\lnot !B}{}}{%
    \indcase{!A}{\lnot !B}{إذا كان $\Sat{{\OLId{latin-looped}{M}}}{\indfrm}[{\OLId{latin-ordinary}{s}}_{\OLMathNumeral{1}}]$، فإن
      $\Sat/{{\OLId{latin-looped}{M}}}{!B}[{\OLId{latin-ordinary}{s}}_{\OLMathNumeral{1}}]$، ومن ثم بفرضية الاستقراء
      $\Sat/{{\OLId{latin-looped}{M}}}{!B}[{\OLId{latin-ordinary}{s}}_{\OLMathNumeral{2}}]$، وبالتالي $\Sat{{\OLId{latin-looped}{M}}}{\indfrm}[{\OLId{latin-ordinary}{s}}_{\OLMathNumeral{2}}]$.}}}

\tagitem{defAnd}{}{%
  \iftag{probAnd}{%
    \indcase!{!A}{!B \land !C}{}}{%
    \indcase{!A}{!B \land !C}{إذا كان $\Sat{{\OLId{latin-looped}{M}}}{\indfrm}[{\OLId{latin-ordinary}{s}}_{\OLMathNumeral{1}}]$،
      فإن $\Sat{{\OLId{latin-looped}{M}}}{!B}[{\OLId{latin-ordinary}{s}}_{\OLMathNumeral{1}}]$ و$\Sat{{\OLId{latin-looped}{M}}}{!C}[{\OLId{latin-ordinary}{s}}_{\OLMathNumeral{1}}]$، ومن ثم بفرضية
      الاستقراء $\Sat{{\OLId{latin-looped}{M}}}{!B}[{\OLId{latin-ordinary}{s}}_{\OLMathNumeral{2}}]$ و$\Sat{{\OLId{latin-looped}{M}}}{!C}[{\OLId{latin-ordinary}{s}}_{\OLMathNumeral{2}}]$. وبالتالي
      $\Sat{{\OLId{latin-looped}{M}}}{\indfrm}[{\OLId{latin-ordinary}{s}}_{\OLMathNumeral{2}}]$.}}}

\tagitem{defOr}{}{%
  \iftag{probOr}{%
    \indcase!{!A}{!B \lor !C}{}}{%
    \indcase{!A}{!B \lor !C}{إذا كان $\Sat{{\OLId{latin-looped}{M}}}{\indfrm}[{\OLId{latin-ordinary}{s}}_{\OLMathNumeral{1}}]$،
      فإن $\Sat{{\OLId{latin-looped}{M}}}{!B}[{\OLId{latin-ordinary}{s}}_{\OLMathNumeral{1}}]$ أو $\Sat{{\OLId{latin-looped}{M}}}{!C}[{\OLId{latin-ordinary}{s}}_{\OLMathNumeral{1}}]$. وبفرضية
      الاستقراء $\Sat{{\OLId{latin-looped}{M}}}{!B}[{\OLId{latin-ordinary}{s}}_{\OLMathNumeral{2}}]$ أو $\Sat{{\OLId{latin-looped}{M}}}{!C}[{\OLId{latin-ordinary}{s}}_{\OLMathNumeral{2}}]$، ومن ثم
      $\Sat{{\OLId{latin-looped}{M}}}{\indfrm}[{\OLId{latin-ordinary}{s}}_{\OLMathNumeral{2}}]$.}}}

\tagitem{defIf}{}{%
  \iftag{probIf}{%
    \indcase!{!A}{!B \lif !C}{}}{%
    \indcase{!A}{!B \lif !C}{إذا كان $\Sat{{\OLId{latin-looped}{M}}}{\indfrm}[{\OLId{latin-ordinary}{s}}_{\OLMathNumeral{1}}]$،
    فإن $\Sat/{{\OLId{latin-looped}{M}}}{!B}[{\OLId{latin-ordinary}{s}}_{\OLMathNumeral{1}}]$ أو $\Sat{{\OLId{latin-looped}{M}}}{!C}[{\OLId{latin-ordinary}{s}}_{\OLMathNumeral{1}}]$. وبفرضية
    الاستقراء $\Sat/{{\OLId{latin-looped}{M}}}{!B}[{\OLId{latin-ordinary}{s}}_{\OLMathNumeral{2}}]$ أو $\Sat{{\OLId{latin-looped}{M}}}{!C}[{\OLId{latin-ordinary}{s}}_{\OLMathNumeral{2}}]$، ومن ثم
    $\Sat{{\OLId{latin-looped}{M}}}{\indfrm}[{\OLId{latin-ordinary}{s}}_{\OLMathNumeral{2}}]$.}}}

\tagitem{defIff}{}{%
  \iftag{probIff}{%
    \indcase!{!A}{!B \liff !C}{}}{%
    \indcase{!A}{!B \liff !C}{إذا كان $\Sat{{\OLId{latin-looped}{M}}}{\indfrm}[{\OLId{latin-ordinary}{s}}_{\OLMathNumeral{1}}]$،
      فإما أن $\Sat{{\OLId{latin-looped}{M}}}{!B}[{\OLId{latin-ordinary}{s}}_{\OLMathNumeral{1}}]$ و$\Sat{{\OLId{latin-looped}{M}}}{!C}[{\OLId{latin-ordinary}{s}}_{\OLMathNumeral{1}}]$، وإما أن
      $\Sat/{{\OLId{latin-looped}{M}}}{!B}[{\OLId{latin-ordinary}{s}}_{\OLMathNumeral{1}}]$ و$\Sat/{{\OLId{latin-looped}{M}}}{!C}[{\OLId{latin-ordinary}{s}}_{\OLMathNumeral{1}}]$. وبفرضية الاستقراء،
      إما أن $\Sat{{\OLId{latin-looped}{M}}}{!B}[{\OLId{latin-ordinary}{s}}_{\OLMathNumeral{2}}]$ و$\Sat{{\OLId{latin-looped}{M}}}{!C}[{\OLId{latin-ordinary}{s}}_{\OLMathNumeral{2}}]$، وإما أن
      $\Sat/{{\OLId{latin-looped}{M}}}{!B}[{\OLId{latin-ordinary}{s}}_{\OLMathNumeral{2}}]$ و$\Sat/{{\OLId{latin-looped}{M}}}{!C}[{\OLId{latin-ordinary}{s}}_{\OLMathNumeral{2}}]$. وفي الحالتين،
      $\Sat{{\OLId{latin-looped}{M}}}{\indfrm}[{\OLId{latin-ordinary}{s}}_{\OLMathNumeral{2}}]$.}}}

\tagitem{defEx}{}{%
  \iftag{probEx}{%
    \indcase!{!A}{\lexists[x][!B]}{}}{%
    \indcase{!A}{\lexists[x][!B]}{إذا كان
      $\Sat{{\OLId{latin-looped}{M}}}{\indfrm}[{\OLId{latin-ordinary}{s}}_{\OLMathNumeral{1}}]$، وُجد ${\OLId{latin-ordinary}{m}} \in \Domain{{\OLId{latin-looped}{M}}}$ بحيث
      $\Sat{{\OLId{latin-looped}{M}}}{!B}[\Subst{{\OLId{latin-ordinary}{s}}_{\OLMathNumeral{1}}}{{\OLId{latin-ordinary}{m}}}{{\OLId{latin-ordinary}{x}}}]$. ليكن
      ${\OLId{latin-ordinary}{s}}_{\OLMathNumeral{1}}' = \Subst{{\OLId{latin-ordinary}{s}}_{\OLMathNumeral{1}}}{{\OLId{latin-ordinary}{m}}}{{\OLId{latin-ordinary}{x}}}$ و${\OLId{latin-ordinary}{s}}_{\OLMathNumeral{2}}' =
      \Subst{{\OLId{latin-ordinary}{s}}_{\OLMathNumeral{2}}}{{\OLId{latin-ordinary}{m}}}{{\OLId{latin-ordinary}{x}}}$. تكون المتغيرات الحرة في~$!B$ من بين
      ${\OLId{latin-ordinary}{x}}_{\OLMathNumeral{1}}$، و\dots، ${\OLId{latin-ordinary}{x}}_{\OLId{latin-ordinary}{n}}$، و${\OLId{latin-ordinary}{x}}$. ولدينا
      ${\OLId{latin-ordinary}{s}}_{\OLMathNumeral{1}}'({\OLId{latin-ordinary}{x}}_{\OLId{latin-ordinary}{i}}) = {\OLId{latin-ordinary}{s}}_{\OLMathNumeral{2}}'({\OLId{latin-ordinary}{x}}_{\OLId{latin-ordinary}{i}})$، لأن ${\OLId{latin-ordinary}{s}}_{\OLMathNumeral{1}}'$ و${\OLId{latin-ordinary}{s}}_{\OLMathNumeral{2}}'$ مغايران بالنسبة
      إلى~${\OLId{latin-ordinary}{x}}$ لـ~${\OLId{latin-ordinary}{s}}_{\OLMathNumeral{1}}$ و${\OLId{latin-ordinary}{s}}_{\OLMathNumeral{2}}$، على التوالي، ولأن
      ${\OLId{latin-ordinary}{s}}_{\OLMathNumeral{1}}({\OLId{latin-ordinary}{x}}_{\OLId{latin-ordinary}{i}}) = {\OLId{latin-ordinary}{s}}_{\OLMathNumeral{2}}({\OLId{latin-ordinary}{x}}_{\OLId{latin-ordinary}{i}})$ بالفرض. ولدينا
      ${\OLId{latin-ordinary}{s}}_{\OLMathNumeral{1}}'({\OLId{latin-ordinary}{x}}) = {\OLId{latin-ordinary}{s}}_{\OLMathNumeral{2}}'({\OLId{latin-ordinary}{x}}) = {\OLId{latin-ordinary}{m}}$ بحسب تعريفنا لـ~${\OLId{latin-ordinary}{s}}_{\OLMathNumeral{1}}'$ و${\OLId{latin-ordinary}{s}}_{\OLMathNumeral{2}}'$.
      وعندئذ تنطبق فرضية الاستقراء على~$!B$ وعلى ${\OLId{latin-ordinary}{s}}_{\OLMathNumeral{1}}'$ و${\OLId{latin-ordinary}{s}}_{\OLMathNumeral{2}}'$،
      فيكون $\Sat{{\OLId{latin-looped}{M}}}{!B}[{\OLId{latin-ordinary}{s}}_{\OLMathNumeral{2}}']$. ومن ثم، لما كان
      ${\OLId{latin-ordinary}{s}}_{\OLMathNumeral{2}}' = \Subst{{\OLId{latin-ordinary}{s}}_{\OLMathNumeral{2}}}{{\OLId{latin-ordinary}{m}}}{{\OLId{latin-ordinary}{x}}}$، يوجد ${\OLId{latin-ordinary}{m}} \in \Domain{{\OLId{latin-looped}{M}}}$ بحيث
      $\Sat{{\OLId{latin-looped}{M}}}{!B}[\Subst{{\OLId{latin-ordinary}{s}}_{\OLMathNumeral{2}}}{{\OLId{latin-ordinary}{m}}}{{\OLId{latin-ordinary}{x}}}]$، ولذلك
      $\Sat{{\OLId{latin-looped}{M}}}{\indfrm}[{\OLId{latin-ordinary}{s}}_{\OLMathNumeral{2}}]$.}}}

\tagitem{defAll}{}{%
  \iftag{probAll}{%
    \indcase!{!A}{\lforall[x][!B]}{}}{%
    \indcase{!A}{\lforall[x][!B]}{إذا كان
      $\Sat{{\OLId{latin-looped}{M}}}{\indfrm}[{\OLId{latin-ordinary}{s}}_{\OLMathNumeral{1}}]$، فإن لكل ${\OLId{latin-ordinary}{m}} \in \Domain{{\OLId{latin-looped}{M}}}$،
      $\Sat{{\OLId{latin-looped}{M}}}{!B}[\Subst{{\OLId{latin-ordinary}{s}}_{\OLMathNumeral{1}}}{{\OLId{latin-ordinary}{m}}}{{\OLId{latin-ordinary}{x}}}]$. ونريد أن نبيّن أنه أيضًا،
      لكل ${\OLId{latin-ordinary}{m}} \in \Domain{{\OLId{latin-looped}{M}}}$، $\Sat{{\OLId{latin-looped}{M}}}{!B}[\Subst{{\OLId{latin-ordinary}{s}}_{\OLMathNumeral{2}}}{{\OLId{latin-ordinary}{m}}}{{\OLId{latin-ordinary}{x}}}]$.
      فليكن ${\OLId{latin-ordinary}{m}} \in \Domain{{\OLId{latin-looped}{M}}}$ اعتباطيًا،
      وانظر في ${\OLId{latin-ordinary}{s}}_{\OLMathNumeral{1}}' = \Subst{{\OLId{latin-ordinary}{s}}_{\OLMathNumeral{1}}}{{\OLId{latin-ordinary}{m}}}{{\OLId{latin-ordinary}{x}}}$ و${\OLId{latin-ordinary}{s}}_{\OLMathNumeral{2}}' =
      \Subst{{\OLId{latin-ordinary}{s}}_{\OLMathNumeral{2}}}{{\OLId{latin-ordinary}{m}}}{{\OLId{latin-ordinary}{x}}}$. لدينا $\Sat{{\OLId{latin-looped}{M}}}{!B}[{\OLId{latin-ordinary}{s}}_{\OLMathNumeral{1}}']$. وتكون
      المتغيرات الحرة في~$!B$ من بين ${\OLId{latin-ordinary}{x}}_{\OLMathNumeral{1}}$، و\dots، ${\OLId{latin-ordinary}{x}}_{\OLId{latin-ordinary}{n}}$، و${\OLId{latin-ordinary}{x}}$.
      ولدينا ${\OLId{latin-ordinary}{s}}_{\OLMathNumeral{1}}'({\OLId{latin-ordinary}{x}}_{\OLId{latin-ordinary}{i}}) = {\OLId{latin-ordinary}{s}}_{\OLMathNumeral{2}}'({\OLId{latin-ordinary}{x}}_{\OLId{latin-ordinary}{i}})$، لأن ${\OLId{latin-ordinary}{s}}_{\OLMathNumeral{1}}'$ و${\OLId{latin-ordinary}{s}}_{\OLMathNumeral{2}}'$ مغايران
      بالنسبة إلى~${\OLId{latin-ordinary}{x}}$ لـ~${\OLId{latin-ordinary}{s}}_{\OLMathNumeral{1}}$ و${\OLId{latin-ordinary}{s}}_{\OLMathNumeral{2}}$، على التوالي، ولأن
      ${\OLId{latin-ordinary}{s}}_{\OLMathNumeral{1}}({\OLId{latin-ordinary}{x}}_{\OLId{latin-ordinary}{i}}) = {\OLId{latin-ordinary}{s}}_{\OLMathNumeral{2}}({\OLId{latin-ordinary}{x}}_{\OLId{latin-ordinary}{i}})$ بالفرض. كما أن
      ${\OLId{latin-ordinary}{s}}_{\OLMathNumeral{1}}'({\OLId{latin-ordinary}{x}}) = {\OLId{latin-ordinary}{s}}_{\OLMathNumeral{2}}'({\OLId{latin-ordinary}{x}}) = {\OLId{latin-ordinary}{m}}$ بحسب تعريفنا لـ~${\OLId{latin-ordinary}{s}}_{\OLMathNumeral{1}}'$ و${\OLId{latin-ordinary}{s}}_{\OLMathNumeral{2}}'$.
      وعندئذ تنطبق فرضية الاستقراء على~$!B$ وعلى ${\OLId{latin-ordinary}{s}}_{\OLMathNumeral{1}}'$ و${\OLId{latin-ordinary}{s}}_{\OLMathNumeral{2}}'$،
      فيكون $\Sat{{\OLId{latin-looped}{M}}}{!B}[{\OLId{latin-ordinary}{s}}_{\OLMathNumeral{2}}']$. وينطبق هذا على كل
      ${\OLId{latin-ordinary}{m}} \in \Domain{{\OLId{latin-looped}{M}}}$، أي إن
      $\Sat{{\OLId{latin-looped}{M}}}{!B}[\Subst{{\OLId{latin-ordinary}{s}}_{\OLMathNumeral{2}}}{{\OLId{latin-ordinary}{m}}}{{\OLId{latin-ordinary}{x}}}]$ لكل~${\OLId{latin-ordinary}{m}} \in
      \Domain{{\OLId{latin-looped}{M}}}$، ومن ثم $\Sat{{\OLId{latin-looped}{M}}}{\indfrm}[{\OLId{latin-ordinary}{s}}_{\OLMathNumeral{2}}]$.}}}
\end{enumerate}
فبمبدأ الاستقراء ثبت أن $\Sat{{\OLId{latin-looped}{M}}}{!A}[{\OLId{latin-ordinary}{s}}_{\OLMathNumeral{1}}]$ إذا وفقط إذا كان
$\Sat{{\OLId{latin-looped}{M}}}{!A}[{\OLId{latin-ordinary}{s}}_{\OLMathNumeral{2}}]$ متى كانت ال!!{variable}s الحرة في~$!A$ من بين
${\OLId{latin-ordinary}{x}}_{\OLMathNumeral{1}}$، و\dots، ${\OLId{latin-ordinary}{x}}_{\OLId{latin-ordinary}{n}}$، وكان ${\OLId{latin-ordinary}{s}}_{\OLMathNumeral{1}}({\OLId{latin-ordinary}{x}}_{\OLId{latin-ordinary}{i}})={\OLId{latin-ordinary}{s}}_{\OLMathNumeral{2}}({\OLId{latin-ordinary}{x}}_{\OLId{latin-ordinary}{i}})$ عندما ${\OLId{latin-ordinary}{i}} = {\OLMathNumeral{1}}$،
و\dots،~${\OLId{latin-ordinary}{n}}$.
\end{proof}

\begin{probtag}{probNot,probOr,probAnd,probIf,probIff,probEx,probAll}
أكمل برهان \olref[fol][syn][ass]{prop:satindep}.
\end{probtag}

\begin{explain}
وأما ال!!^{sentence}s فخالية من المتغيرات الحرة؛ لذلك يتفق كل
إسنادين على جميع متغيراتها الحرة، إذ لا شيء منها أصلًا. فتقتضي
القضية المتقدمة أن إرضاء !!a{sentence} في بنية، أو عدم إرضائها،
لا يتأثر باختيار إسناد المتغيرات. نثبت هذه النتيجة فيما يأتي،
وبها يصح تعريف إرضاء !!a{sentence} في !!a{structure}
من غير ذكر الإسناد.
\end{explain}

\begin{cor}
\ollabel{cor:sat-sentence}
إذا كانت $!A$ !!a{sentence} وكان ${\OLId{latin-ordinary}{s}}$ إسنادًا للمتغيرات، فإن
$\Sat{{\OLId{latin-looped}{M}}}{!A}[{\OLId{latin-ordinary}{s}}]$ إذا وفقط إذا كان $\Sat{{\OLId{latin-looped}{M}}}{!A}[{\OLId{latin-ordinary}{s}}']$ لكل إسناد
للمتغيرات~${\OLId{latin-ordinary}{s}}'$.
\end{cor}

\begin{proof}
  لنأخذ ${\OLId{latin-ordinary}{s}}'$ إسنادًا اعتباطيًا. ولما كانت $!A$ !!a{sentence}،
  لم يكن فيها متغير حر؛ فيتفق كل إسناد~${\OLId{latin-ordinary}{s}}'$ مع الإسناد~${\OLId{latin-ordinary}{s}}$
  على جميع المتغيرات الحرة في~$!A$ اتفاقًا خاليًا من الشروط.
  فقد تحقق فرض \olref{prop:satindep}، ومنه
  $\Sat{{\OLId{latin-looped}{M}}}{!A}[{\OLId{latin-ordinary}{s}}]$ إذا وفقط إذا كان $\Sat{{\OLId{latin-looped}{M}}}{!A}[{\OLId{latin-ordinary}{s}}']$.
\end{proof}

\begin{defn}
\ollabel{defn:satisfaction}
إذا كانت $!A$ !!a{sentence}، قلنا إن !!a{structure}~$\Struct M$
\emph{ترضي}~$!A$، $\Sat{{\OLId{latin-looped}{M}}}{!A}$، إذا وفقط إذا كان
$\Sat{{\OLId{latin-looped}{M}}}{!A}[{\OLId{latin-ordinary}{s}}]$ لجميع إسنادات المتغيرات~${\OLId{latin-ordinary}{s}}$.
\end{defn}

إذا كان $\Sat{{\OLId{latin-looped}{M}}}{!A}$، قلنا أيضًا ببساطة إن
\emph{$!A$ صادقة في~$\Struct{M}$.} ونمدّ هذا التعريف
من ال!!{sentence}s المفردة إلى مجموعات ال!!{sentence}s.

\begin{defn}
\ollabel{defn:sat}
إذا كانت ${\OLId{greek-variable}{Gamma}}$ مجموعة من ال!!{sentence}s، قلنا إن
!!a{structure}~$\Struct M$ \emph{ترضي}~${\OLId{greek-variable}{Gamma}}$،
$\Sat{{\OLId{latin-looped}{M}}}{{\OLId{greek-variable}{Gamma}}}$، إذا وفقط إذا كان $\Sat{{\OLId{latin-looped}{M}}}{!A}$ لكل
$!A \in {\OLId{greek-variable}{Gamma}}$.
\end{defn}

\begin{prop}\ollabel{prop:sentence-sat-true}
  لتكن $\Struct{M}$ !!a{structure}، ولتكن $!A$ !!a{sentence}، وليكن
  ${\OLId{latin-ordinary}{s}}$ إسنادًا للمتغيرات. يكون $\Sat{{\OLId{latin-looped}{M}}}{!A}$ إذا وفقط إذا كان
  $\Sat{{\OLId{latin-looped}{M}}}{!A}[{\OLId{latin-ordinary}{s}}]$.
\end{prop}

\begin{proof}
تمرين.
\end{proof}

\begin{prob}
أثبت \olref[fol][syn][ass]{prop:sentence-sat-true}.
\end{prob}

\begin{prop}\ollabel{prop:sat-quant}
  افترض أن $!A({\OLId{latin-ordinary}{x}})$ لا تحتوي حرًا إلا على~${\OLId{latin-ordinary}{x}}$، وأن $\Struct M$
  !!a{structure}. عندئذ:
  \begin{tagenumerate}{prvEx,prvAll}
    \tagitem{prvEx}{$\Sat{{\OLId{latin-looped}{M}}}{\lexists[{\OLId{latin-ordinary}{x}}][!A({\OLId{latin-ordinary}{x}})]}$ إذا وفقط إذا كان
      $\Sat{{\OLId{latin-looped}{M}}}{!A({\OLId{latin-ordinary}{x}})}[{\OLId{latin-ordinary}{s}}]$ لإسناد متغيرات واحد على الأقل~${\OLId{latin-ordinary}{s}}$.}{}
    \tagitem{prvAll}{$\Sat{{\OLId{latin-looped}{M}}}{\lforall[{\OLId{latin-ordinary}{x}}][!A({\OLId{latin-ordinary}{x}})]}$ إذا وفقط إذا كان
      $\Sat{{\OLId{latin-looped}{M}}}{!A({\OLId{latin-ordinary}{x}})}[{\OLId{latin-ordinary}{s}}]$ لجميع إسنادات المتغيرات~${\OLId{latin-ordinary}{s}}$.}{}
\end{tagenumerate}
\end{prop}

\begin{proof}
تمرين.
\end{proof}

\begin{prob}
أثبت \olref[fol][syn][ass]{prop:sat-quant}.
\end{prob}

\begin{prob}
\DeclareRobustCommand{\VDash}{\mathrel{||}\joinrel\Relbar}
لتكن $\Lang L$ لغة ليس فيها !!{function}s. ولتكن معطاة
!!a{structure}~$\Struct M$، و!!a{constant}~${\OLId{latin-ordinary}{c}}$، و
${\OLId{latin-ordinary}{a}} \in \Domain {\OLId{latin-looped}{M}}$، فلتكن $\Struct M[{\OLId{latin-ordinary}{a}}/{\OLId{latin-ordinary}{c}}]$ ال!!{structure}
الموافقة لـ~$\Struct M$ في كل شيء، سوى أن $\Assign{{\OLId{latin-ordinary}{c}}}{{\OLId{latin-looped}{M}}[{\OLId{latin-ordinary}{a}}/{\OLId{latin-ordinary}{c}}]} = {\OLId{latin-ordinary}{a}}$. وعرّف
$\Struct M \VDash !A$ لل!!{sentence}s~$!A$ بما يأتي:
\begin{enumerate}
\tagitem{prvFalse}{%
  \indcase{!A}{\lfalse}{ليس $\Struct{M} \VDash \indfrm$.}}{}

\tagitem{prvTrue}{%
  \indcase{!A}{\ltrue}{$\Struct{M} \VDash \indfrm$.}}{}

\item \indcase{!A}{\Atom{R}{d_1, \dots, d_n}}{$\Struct M \VDash \indfrm$
  إذا وفقط إذا كان $\langle \Assign{{\OLId{latin-ordinary}{d}}_{\OLMathNumeral{1}}}{{\OLId{latin-looped}{M}}}, \dots,
  \Assign{{\OLId{latin-ordinary}{d}}_{\OLId{latin-ordinary}{n}}}{{\OLId{latin-looped}{M}}} \rangle \in \Assign{{\OLId{latin-looped}{R}}}{{\OLId{latin-looped}{M}}}$.}

\item \indcase{!A}{\eq[d_1][d_2]}{$\Struct M \VDash \indfrm$ إذا
  وفقط إذا كان $\Assign{{\OLId{latin-ordinary}{d}}_{\OLMathNumeral{1}}}{{\OLId{latin-looped}{M}}} = \Assign{{\OLId{latin-ordinary}{d}}_{\OLMathNumeral{2}}}{{\OLId{latin-looped}{M}}}$.}

\tagitem{prvNot}{%
  \indcase{!A}{\lnot !B}{$\Struct{M} \VDash \indfrm$ إذا وفقط إذا
    لم يكن $\Struct{M} \VDash {!B}$.}}{}

\tagitem{prvAnd}{%
  \indcase{!A}{(!B \land !C)}{$\Struct{M} \VDash
    \indfrm$ إذا وفقط إذا كان $\Struct{M} \VDash!B$ و
    $\Struct{M} \VDash !C$.}}{}

\tagitem{prvOr}{%
  \indcase{!A}{(!B \lor !C)}{$\Struct{M} \VDash \indfrm$ إذا وفقط
    إذا كان $\Struct{M} \VDash !B$ أو $\Struct{M} \VDash !C$
    (أو كلاهما).}}{}

\tagitem{prvIf}{%
  \indcase{!A}{(!B \lif !C)}{$\Struct{M} \VDash \indfrm$ إذا
    وفقط إذا لم يكن $\Struct {M} \VDash !B$، أو كان
    $\Struct M \VDash !C$ (أو كلاهما).}}{}

\tagitem{prvIff}{%
  \indcase{!A}{(!B \liff !C)}{$\Struct{M} \VDash {\indfrm}$ إذا
    وفقط إذا كان إما كل من $\Struct{M} \VDash {!B}$ و
    $\Struct{M} \VDash {!C}$، وإما لا $\Struct{M} \VDash {!B}$
    ولا $\Struct{M} \VDash {!C}$.}}{}

\tagitem{prvAll}{%
  \indcase{!A}{\lforall[x][!B]}{$\Struct{M} \VDash {\indfrm}$ إذا
    وفقط إذا كان، لكل ${\OLId{latin-ordinary}{a}} \in \Domain{{\OLId{latin-looped}{M}}}$،
    $\Struct{M[a/c]} \VDash \Subst{!B}{{\OLId{latin-ordinary}{c}}}{{\OLId{latin-ordinary}{x}}}$، إذا لم يرد~${\OLId{latin-ordinary}{c}}$
    في~$!B$.}}{}

\tagitem{prvEx}{%
  \indcase{!A}{\lexists[x][!B]}{$\Struct{M} \VDash
  {\indfrm}$ إذا وفقط إذا وُجد ${\OLId{latin-ordinary}{a}} \in \Domain {\OLId{latin-looped}{M}}$ بحيث
  $\Struct{M[a/c]} \VDash \Subst{!B}{{\OLId{latin-ordinary}{c}}}{{\OLId{latin-ordinary}{x}}}$، إذا لم يرد~${\OLId{latin-ordinary}{c}}$
  في~$!B$.}}{}
\end{enumerate}
لتكن ${\OLId{latin-ordinary}{x}}_{\OLMathNumeral{1}}$، و\dots، ${\OLId{latin-ordinary}{x}}_{\OLId{latin-ordinary}{n}}$ جميع ال!!{variable}s الحرة في~$!A$،
ولتكن ${\OLId{latin-ordinary}{c}}_{\OLMathNumeral{1}}$، و\dots، ${\OLId{latin-ordinary}{c}}_{\OLId{latin-ordinary}{n}}$ رموزًا ثابتة لا ترد في~$!A$، ولتكن
${\OLId{latin-ordinary}{a}}_{\OLMathNumeral{1}}$، و\dots، ${\OLId{latin-ordinary}{a}}_{\OLId{latin-ordinary}{n}} \in \Domain {\OLId{latin-looped}{M}}$، وليكن ${\OLId{latin-ordinary}{s}}({\OLId{latin-ordinary}{x}}_{\OLId{latin-ordinary}{i}}) = {\OLId{latin-ordinary}{a}}_{\OLId{latin-ordinary}{i}}$.

بيّن أن $\Sat{{\OLId{latin-looped}{M}}}{!A}[{\OLId{latin-ordinary}{s}}]$ إذا وفقط إذا كان
$\Struct M[{\OLId{latin-ordinary}{a}}_{\OLMathNumeral{1}}/{\OLId{latin-ordinary}{c}}_{\OLMathNumeral{1}},\dots,{\OLId{latin-ordinary}{a}}_{\OLId{latin-ordinary}{n}}/{\OLId{latin-ordinary}{c}}_{\OLId{latin-ordinary}{n}}]
\VDash \Subst{\Subst{!A}{{\OLId{latin-ordinary}{c}}_{\OLMathNumeral{1}}}{{\OLId{latin-ordinary}{x}}_{\OLMathNumeral{1}}}\dots}{{\OLId{latin-ordinary}{c}}_{\OLId{latin-ordinary}{n}}}{{\OLId{latin-ordinary}{x}}_{\OLId{latin-ordinary}{n}}}$.

(وبهذه المسألة يتبيّن أن منطق الرتبة الأولى يقبل دلالة لا تحتاج
إلى إسنادات المتغيرات.)
\end{prob}

\begin{prob}
افترض أن ${\OLId{latin-ordinary}{f}}$ رمز دالة لا يرد في~$!A({\OLId{latin-ordinary}{x}},{\OLId{latin-ordinary}{y}})$. بيّن أنه توجد
!!a{structure}~$\Struct{M}$ بحيث
$\Sat{{\OLId{latin-looped}{M}}}{\lforall[{\OLId{latin-ordinary}{x}}][\lexists[{\OLId{latin-ordinary}{y}}][!A({\OLId{latin-ordinary}{x}},{\OLId{latin-ordinary}{y}})]]}$ إذا وفقط إذا وُجدت
~$\Struct M'$ بحيث $\Sat{{\OLId{latin-looped}{M}}'}{\lforall[{\OLId{latin-ordinary}{x}}][!A({\OLId{latin-ordinary}{x}},{\OLId{latin-ordinary}{f}}({\OLId{latin-ordinary}{x}}))]}$.

(وهذه حالة خاصة من المبرهنة المعروفة باسم مبرهنة سكولم؛ وتسمى
$\lforall[{\OLId{latin-ordinary}{x}}][!A({\OLId{latin-ordinary}{x}},{\OLId{latin-ordinary}{f}}({\OLId{latin-ordinary}{x}}))]$ \emph{صيغة سكولم السوية}
لـ~$\lforall[{\OLId{latin-ordinary}{x}}][\lexists[{\OLId{latin-ordinary}{y}}][!A({\OLId{latin-ordinary}{x}},{\OLId{latin-ordinary}{y}})]]$.)
\end{prob}

\end{document}
